% \subsection{Bối cảnh nghiên cứu}

% Phân đoạn ngữ nghĩa là bài toán gán nhãn cho từng pixel trong ảnh, đòi hỏi mô hình đồng thời nắm bắt ngữ cảnh toàn cục và bảo toàn đặc trưng cục bộ. Hiện nay, bài toán này được tiếp cận bởi nhiều họ kiến trúc như CNN-based, Transformer-based,.... Tuy nhiên, trong điều kiện có ràng buộc về độ trễ và tài nguyên tính toán, các kiến trúc CNN-based hoặc hybrid nhẹ vẫn là lựa chọn thực tiễn nhờ hiệu quả suy luận cao.

% Đối với bài toán phân đoạn fine-grained, thách thức không chỉ nằm ở việc mô hình hóa ngữ nghĩa mà còn ở khả năng duy trì các đặc trưng chi tiết như biên và texture. Trong khi đó, quá trình trừu tượng hóa và giảm độ phân giải đặc trưng trong các mô hình nhẹ thường làm suy giảm thông tin tần số cao, dẫn đến mất các tín hiệu phân biệt quan trọng ở mức cục bộ.

% Phân đoạn ảnh thực phẩm là một trường hợp điển hình của bài toán này. So với các bối cảnh phân đoạn thông thường, ảnh món ăn có mức độ chồng lấp cao, ranh giới không rõ ràng, biến thiên hình thái lớn trong cùng lớp, trong khi nhiều lớp khác nhau lại có độ tương đồng thị giác cao. Vì vậy, khả năng khai thác texture và các biến thiên cục bộ đóng vai trò quan trọng trong việc phân biệt các thành phần nguyên liệu \cite{zhuang2020transfer}.

% Tập dữ liệu \textbf{FoodSeg103} là một benchmark tiêu biểu cho phân đoạn ảnh thực phẩm ở mức ingredient với số lượng lớp lớn và tính chất fine-grained rõ rệt. Đặc điểm của tập dữ liệu này khiến nó trở thành tiêu chuẩn phù hợp để đánh giá mức độ bảo toàn đặc trưng cục bộ của các mô hình phân đoạn.

% Về mặt ứng dụng thực tiễn, nhu cầu tự động hóa trong xử lý, giám sát và nhận diện thực phẩm đặt ra yêu cầu đối với các mô hình vừa chính xác vừa có khả năng triển khai thực tế \cite{mahalik2010trends, stemmer2018vision, gao2019smartfridge, wang2019smartrefrigerator}. Tuy nhiên, các mô hình hiệu năng cao thường đi kèm chi phí tính toán lớn, trong khi các mô hình lightweight lại dễ suy giảm khả năng biểu diễn texture do giới hạn kiến trúc \cite{he2015residual, dong2020mobilenetv2}. Khoảng trống nghiên cứu vì vậy nằm ở việc thiết kế các phương pháp phân đoạn nhẹ nhưng vẫn bảo toàn và tăng cường được thông tin texture cần thiết cho bài toán phân đoạn ảnh món ăn fine-grained.

% % ------------------------------
% Phân đoạn ngữ nghĩa (semantic segmentation) là một bài toán quan trọng trong lĩnh vực thị giác máy tính, trong đó mục tiêu là phân lớp cho từng pixel trong ảnh.

% Dân số toàn cầu dự kiến đạt 9,1 tỷ người vào năm 2050, đòi hỏi sản lượng lương thực tăng khoảng 70\% \cite{godfray2010foodsecurity, alexandratos1995world}. Vì vậy, việc tự động hóa chuỗi cung ứng thực phẩm --- từ thu hoạch, phân loại đến kiểm soát chất lượng --- trở nên cần thiết \cite{mahalik2010trends, stemmer2018vision}. Trong đó, phân vùng ngữ nghĩa giúp hệ thống robot và camera nhận diện chính xác thực phẩm, giảm sai sót trong sản xuất, đồng thời hỗ trợ các thiết bị IoT như tủ lạnh thông minh theo dõi độ tươi và dinh dưỡng theo thời gian thực \cite{gao2019smartfridge, wang2019smartrefrigerator}.

% Từ góc độ thuật toán, hình ảnh thực phẩm có độ phức tạp cao do các đối tượng thường chồng lấp, đa dạng hình thái trong cùng nhóm và có nhiều đặc điểm tương đồng giữa các nhóm khác nhau \cite{zhuang2020transfer}. Mặc dù các CNN sâu có khả năng trích xuất đặc trưng mạnh, việc triển khai trong hệ thống thực tế đòi hỏi cân bằng giữa độ chính xác, tốc độ xử lý và tài nguyên phần cứng \cite{he2015residual}. Vì vậy, nghiên cứu các mô hình học sâu gọn nhẹ (lightweight models) nhằm giảm chi phí tính toán nhưng vẫn đảm bảo hiệu năng nhận diện là một hướng nghiên cứu quan trọng \cite{dong2020mobilenetv2}.

% Tập dữ liệu \textbf{FoodSeg103} \cite{benchmarkfood} được xây dựng để hỗ trợ nghiên cứu phân đoạn thực phẩm ở mức pixel với \textbf{103 loại nguyên liệu khác nhau}. Tuy nhiên, bài toán này khá khó vì hình ảnh thực phẩm có đặc trưng thị giác phức tạp:
% \begin{itemize}
%     \item Các nguyên liệu thường chồng lấp lên nhau;
%     \item Một nguyên liệu có thể có nhiều hình thái khác nhau;
%     \item Hai nguyên liệu khác nhau đôi khi có hình dạng rất giống nhau.
% \end{itemize}

% Các mô hình phân đoạn hiện đại như \textbf{DeepLabV3} đạt độ chính xác cao nhờ sử dụng các kỹ thuật như atrous convolution và mô-đun ASPP để khai thác thông tin đa tỉ lệ. Tuy nhiên, các mô hình này thường sử dụng backbone nặng (như ResNet), dẫn đến chi phí tính toán lớn và khó triển khai trên thiết bị di động. Do đó, việc nghiên cứu các backbone lightweight nhằm giảm chi phí tính toán trong khi vẫn giữ được độ chính xác là một hướng nghiên cứu quan trọng.

% \subsection{Vấn đề nghiên cứu}

% Mặc dù DeepLabV3 đạt hiệu năng tốt trong nhiều bài toán phân đoạn ngữ nghĩa, mô hình này thường sử dụng các backbone lớn, khiến thời gian suy luận và kích thước mô hình khá cao. Điều này gây khó khăn khi triển khai trên các thiết bị có tài nguyên hạn chế như điện thoại thông minh hoặc thiết bị nhúng.

% Cần nghiên cứu xem liệu việc thay thế backbone của DeepLabV3 bằng các kiến trúc nhẹ hơn (như MobileNet) có thể cải thiện tốc độ suy luận mà vẫn giữ được độ chính xác chấp nhận được hay không. Ngoài ra, việc thay đổi độ phân giải đầu vào của ảnh cũng có thể ảnh hưởng đến sự cân bằng giữa tốc độ và độ chính xác của mô hình, đặc biệt khi triển khai trên thiết bị di động.

\subsection{Bối cảnh nghiên cứu}

Phân đoạn ngữ nghĩa là bài toán gán nhãn cho từng pixel trong ảnh, yêu cầu mô hình đồng thời nắm bắt ngữ cảnh toàn cục và bảo toàn đặc trưng cục bộ. Có nhiều hướng tiếp cận cho bài toán này, từ các mô hình học máy đơn giản, đến các mô hình học sâu CNN-based hay Transformer-based. Trong đó, các mô hình hiệu năng cao thường tập trung vào việc tăng cường khả năng biểu diễn ngữ nghĩa và khai thác ngữ cảnh toàn cục, còn các kiến trúc nhẹ hướng đến mục tiêu giảm chi phí tính toán để phù hợp với các kịch bản triển khai thực tế.

Trong lĩnh vực thực phẩm, nhu cầu tự động hóa trong phân loại, giám sát và kiểm soát chất lượng đã thúc đẩy sự phát triển của các hệ thống thị giác máy có khả năng nhận diện thành phần thực phẩm một cách chính xác và hiệu quả \cite{mahalik2010trends, stemmer2018vision, gao2019smartfridge, wang2019smartrefrigerator}. Phân đoạn ngữ nghĩa trong lĩnh vực này có các thách thức do các thành phần nguyên liệu thường chồng lấp, ranh giới không rõ ràng, hình thái biến đổi mạnh theo cách chế biến và nhiều lớp có độ tương đồng thị giác cao \cite{zhuang2020transfer}. Vì vậy, đây là một bối cảnh điển hình của bài toán phân đoạn fine-grained, trong đó tín hiệu phân biệt giữa các lớp không chỉ nằm ở hình dạng tổng thể mà còn ở các đặc trưng cục bộ như texture.

Tập dữ liệu \textbf{FoodSeg103} \cite{benchmarkfood} được xây dựng như một benchmark cho phân đoạn ảnh thực phẩm theo từng thành phần gồm 103 lớp và nhãn pixel-wise. Với số lượng lớp lớn cùng tính chất fine-grained rõ rệt, phù hợp để khảo sát khả năng của các mô hình phân đoạn trong việc xử lý dữ liệu thực phẩm có độ phức tạp cao.

\subsection{Vấn đề nghiên cứu}

Mặc dù các mô hình phân đoạn nhẹ có ưu thế về tốc độ suy luận và khả năng triển khai trên thiết bị tài nguyên hạn chế, hiệu quả của chúng trên bài toán phân đoạn ảnh thực phẩm fine-grained vẫn chưa được làm rõ đầy đủ. Trên các tập dữ liệu như FoodSeg103, chất lượng phân đoạn không chỉ phụ thuộc vào khả năng nắm bắt ngữ cảnh mà còn phụ thuộc mạnh vào việc bảo toàn các tín hiệu cục bộ như biên và texture, vốn là cơ sở để phân biệt các lớp nguyên liệu có độ tương đồng thị giác cao.

Trong khi đó, các kiến trúc lightweight thường đạt hiệu quả tính toán bằng cách giảm độ phức tạp của backbone hoặc giảm độ phân giải đặc trưng sớm, điều này có thể làm suy giảm khả năng biểu diễn các đặc trưng cục bộ quan trọng \cite{he2015residual, dong2020mobilenetv2}. Vì vậy, vấn đề đặt ra là liệu các hướng tiếp cận phân đoạn ngữ nghĩa nhẹ có thể duy trì chất lượng phân đoạn chấp nhận được trên FoodSeg103 hay không, và ở mức độ nào việc tăng cường biểu diễn texture có thể cải thiện sự đánh đổi giữa độ chính xác và hiệu quả suy luận.


% \subsection{Câu hỏi nghiên cứu}

% \textbf{Câu hỏi chính:} Liệu các backbone lightweight có thể cải thiện tốc độ suy luận của DeepLabV3 trong khi vẫn duy trì hiệu năng phân đoạn chấp nhận được trên tập dữ liệu FoodSeg103 hay không?

\section{Câu hỏi nghiên cứu}
\textbf{Câu hỏi chính:}  
Các mô hình phân đoạn ngữ nghĩa CNN-based nhẹ có thể đạt được sự cân bằng hợp lý giữa hiệu quả tính toán và chất lượng phân đoạn trên tập dữ liệu FoodSeg103 hay không?

% \textbf{Câu hỏi phụ:}
% \begin{enumerate}
%     \item Việc thay thế backbone của DeepLabV3 bằng MobileNetV2 hoặc MobileNetV3 ảnh hưởng như thế nào đến độ chính xác của mô hình?
%     \item Sự đánh đổi giữa tốc độ suy luận và độ chính xác khi sử dụng các backbone lightweight là như thế nào?
%     \item Backbone nào phù hợp nhất cho việc triển khai DeepLabV3 trên thiết bị di động?
% \end{enumerate}

% \subsection{Mục tiêu nghiên cứu}

% \begin{enumerate}
%     \item Triển khai mô hình \textbf{DeepLabV3 baseline} cho bài toán phân đoạn ngữ nghĩa trên tập dữ liệu FoodSeg103.
%     \item Thay thế backbone của DeepLabV3 bằng các kiến trúc lightweight như \textbf{MobileNetV2} và \textbf{MobileNetV3}.
%     \item Thực nghiệm và so sánh hiệu năng của các mô hình theo các tiêu chí độ chính xác và tốc độ suy luận.
%     \item Phân tích ảnh hưởng của độ phân giải đầu vào đến hiệu năng của mô hình.
%     \item Đề xuất cấu hình mô hình phù hợp cho việc triển khai trên thiết bị di động.
% \end{enumerate}

\textbf{Các câu hỏi cụ thể:}
\begin{enumerate}
\item Các kiến trúc CNN-based nhẹ khác nhau ảnh hưởng như thế nào đến chất lượng phân đoạn trên các lớp thực phẩm có tính fine-grained cao?
\item Mức độ suy giảm hiệu năng của các mô hình lightweight có liên quan như thế nào đến việc mất thông tin texture và chi tiết cục bộ?
\item Việc điều chỉnh backbone, độ phân giải đầu vào và cơ chế tăng cường đặc trưng cục bộ có cải thiện được sự đánh đổi giữa độ chính xác và tốc độ suy luận hay không?
\item Cấu hình mô hình nào phù hợp nhất cho bài toán phân đoạn ảnh thực phẩm trong điều kiện tài nguyên tính toán hạn chế?
\end{enumerate}

\section{Mục tiêu nghiên cứu}
\textbf{Mục tiêu tổng quát:}  
Khảo sát và đánh giá các hướng tiếp cận phân đoạn ngữ nghĩa CNN-based nhẹ cho ảnh món ăn fine-grained, đồng thời phân tích vai trò của đặc trưng texture trong sự đánh đổi giữa hiệu năng phân đoạn và hiệu quả tính toán.

\textbf{Các mục tiêu cụ thể:}
\begin{enumerate}
\item Xây dựng quy trình thực nghiệm cho bài toán phân đoạn ngữ nghĩa trên tập dữ liệu FoodSeg103.
\item Lựa chọn và triển khai một số kiến trúc CNN-based nhẹ làm mô hình đối sánh cho bài toán.
\item Đánh giá ảnh hưởng của backbone và độ phân giải đầu vào đến các chỉ số chất lượng phân đoạn và chi phí suy luận.
\item Khảo sát khả năng cải thiện hiệu năng thông qua các cơ chế tăng cường đặc trưng cục bộ theo hướng texture-aware.
\item Đề xuất một cấu hình mô hình khả thi cho các bối cảnh triển khai trên thiết bị có tài nguyên hạn chế.
\end{enumerate}

% \subsection{Tổng quan nghiên cứu liên quan}

% Trong lĩnh vực phân đoạn ngữ nghĩa, nhiều kiến trúc mạng nơ-ron sâu đã được đề xuất như Fully Convolutional Network (FCN), U-Net và DeepLab. DeepLabV3 là một trong những mô hình nổi bật nhờ sử dụng \textbf{atrous convolution} để mở rộng receptive field và mô-đun \textbf{Atrous Spatial Pyramid Pooling (ASPP)} để khai thác thông tin ngữ cảnh đa tỉ lệ.

% Tuy nhiên, các backbone thường dùng trong DeepLabV3 như ResNet có chi phí tính toán cao \cite{he2015residual}. Các nghiên cứu gần đây đã đề xuất sử dụng các kiến trúc lightweight như MobileNet để giảm chi phí tính toán mà vẫn duy trì hiệu năng tốt \cite{dong2020mobilenetv2}. MobileNetV2 và MobileNetV3 được thiết kế đặc biệt cho các thiết bị di động, sử dụng các kỹ thuật như \textbf{depthwise separable convolution} và \textbf{inverted residual block} nhằm giảm số lượng tham số và tăng tốc độ suy luận.

% Ngoài ra, nhiều nghiên cứu ứng dụng đã chứng minh tiềm năng của học sâu trong nhận diện thực phẩm và quản lý lãng phí, từ bài toán nhận diện ảnh món ăn đến các hệ thống nhà bếp thông minh \cite{yanai2015foodimage, lubura2022foodrecognition, chakraborty2024aikitchen}.


\section{Nghiên cứu liên quan}
Các nghiên cứu về phân đoạn ngữ nghĩa hiện đại chủ yếu phát triển theo hai hướng chính. Hướng thứ nhất tập trung vào nâng cao chất lượng biểu diễn ngữ nghĩa thông qua các mô hình CNN-based có khả năng tổng hợp ngữ cảnh đa tỉ lệ. Hướng này thường đạt hiệu năng cao nhưng đi kèm chi phí tính toán lớn. Hướng thứ hai tập trung vào thiết kế các kiến trúc lightweight nhằm giảm số tham số, giảm FLOPs và tăng tốc độ suy luận để phục vụ triển khai thực tế trên thiết bị di động hoặc thiết bị nhúng.

Trong bối cảnh ảnh thực phẩm, các nghiên cứu trước đây cho thấy đây là miền dữ liệu có độ phức tạp cao, nơi sự khác biệt giữa các lớp thường không được xác định rõ bởi hình dạng tổng thể mà bởi các đặc trưng cục bộ như texture, màu sắc và trạng thái bề mặt \cite{zhuang2020transfer}. Điều này khiến các chiến lược giảm độ phân giải đặc trưng quá sớm trong mô hình nhẹ có nguy cơ làm mất thông tin phân biệt quan trọng.

Do đó, một hướng nghiên cứu hợp lý là kết hợp ưu thế về hiệu quả tính toán của kiến trúc lightweight với các cơ chế giúp tăng cường biểu diễn đặc trưng cục bộ. Trong phạm vi nghiên cứu này, các mô hình CNN-based nhẹ được xem như nền tảng thực nghiệm, còn hướng phát triển theo texture-aware được xem như giả thuyết cần kiểm chứng về mặt thực nghiệm trên FoodSeg103.
